\begin{table*}[t]
\centering
\small
\begin{threeparttable}
\caption{Scenario-level gains over CDP for the comprehensive policy library. The 6.20\% mean of relative scenario gains differs from the 3.0\% gain of pooled means.}
\label{tab:scenario_improvement_cdp}
\begin{tabular}{lccccc}
\toprule
\rowcolor{TableHeader}
Coupling pattern 
& Mean EP gain 
& Median EP gain 
& Min. EP gain 
& Max. EP gain 
& Positive cases \\
\midrule
\rowcolor{RowScenario}
DCAC & 5.06\% & 4.60\% & $-9.13\%$ & 19.37\% & 72.22\% \\
\rowcolor{RowBaseline}
DCDC & 5.28\% & 2.69\% & $-8.58\%$ & 27.04\% & 55.56\% \\
\rowcolor{RowProposed}
DDAC & \best{7.97\%} & \best{9.69\%} & $-10.36\%$ & 23.38\% & 72.22\% \\
\rowcolor{RowRisk}
DDDC & 5.41\% & 1.38\% & $-22.46\%$ & \best{34.34\%} & 50.00\% \\
\rowcolor{RowPhysical}
RC & 7.29\% & 3.00\% & $-25.28\%$ & 29.46\% & 72.22\% \\
\midrule
\rowcolor{RowProposed}
Overall & \best{6.20\%} & \best{3.77\%} & $-25.28\%$ & \best{34.34\%} & \best{64.44\%} \\
\bottomrule
\end{tabular}
\begin{tablenotes}
\footnotesize
\item Note: Positive cases are scenario-budget combinations in which the composite EP score is lower than CDP. A positive fraction is descriptive and is not a significance test. Overall values summarize a different aggregation than pooled strategy means.
\end{tablenotes}
\end{threeparttable}
\end{table*}
